%------------------------
% Resume - "clean" variant (2-page)
% Based on Anubhav Singh's MIT-licensed template (github.com/xprilion),
% filled with Swayam Singh's data; header swapped for a centered name +
% FontAwesome contact row (matching the original design).
% Tailored variant: Cloud Infra Engineer (ML Platform).
%------------------------
\documentclass[a4paper,11pt]{article}

\usepackage{latexsym}
\usepackage[empty]{fullpage}
\usepackage{titlesec}
\usepackage[usenames,dvipsnames]{color}
\usepackage{enumitem}
\usepackage[hidelinks]{hyperref}
\usepackage{fancyhdr}
\usepackage{fontawesome5}

\pagestyle{fancy}
\fancyhf{}
\renewcommand{\headrulewidth}{0pt}
\renewcommand{\footrulewidth}{0pt}

% Margins
\addtolength{\oddsidemargin}{-0.530in}
\addtolength{\evensidemargin}{-0.375in}
\addtolength{\textwidth}{1in}
\addtolength{\topmargin}{-.45in}
\addtolength{\textheight}{1in}

\urlstyle{rm}
\raggedbottom
\raggedright
\setlength{\tabcolsep}{0in}

% Section formatting
\titleformat{\section}{\vspace{-6pt}\scshape\raggedright\large}{}{0em}{}[\color{black}\titlerule\vspace{-4pt}]

%------ custom commands ------
\newcommand{\resumeItem}[2]{\item\small{\textbf{#1}: #2}\vspace{-1pt}}
\newcommand{\resumeItemPlain}[1]{\item\small{#1}\vspace{-1pt}}
\newcommand{\resumeSubheading}[4]{%
  \vspace{1pt}\item[]
    \begin{tabular*}{0.99\textwidth}{l@{\extracolsep{\fill}}r}
      \textbf{#1} & #2 \\
      {\small #3} & {\small #4} \\
    \end{tabular*}\vspace{-1pt}}
\newcommand{\resumeSubItem}[2]{\resumeItem{#1}{#2}}
\renewcommand{\labelitemii}{$\circ$}

\newcommand{\resumeHeadingListStart}{\begin{itemize}[leftmargin=0pt, label={}, itemsep=3pt, topsep=3pt, parsep=0pt]}
\newcommand{\resumeHeadingListEnd}{\end{itemize}}
\newcommand{\resumeItemListStart}{\begin{itemize}[leftmargin=1.3em, label=$\bullet$, itemsep=1.5pt, topsep=2pt, parsep=0pt]}
\newcommand{\resumeItemListEnd}{\end{itemize}\vspace{-3pt}}

\begin{document}

%---------- HEADING (centered name + icon contact row) ----------
\begin{center}
  {\Huge\bfseries Swayam Singh}\\[4pt]
  \small
  \faEnvelope~\href{mailto:singhswayam008@gmail.com}{singhswayam008@gmail.com}
  \hspace{1.2em}
  \faPhone*~+91~9116277756
  \hspace{1.2em}
  \href{https://github.com/SwayamInSync}{\faGithub~GitHub}
  \hspace{1.2em}
  \href{https://www.linkedin.com/in/swayam-singh-406610213/}{\faLinkedin~LinkedIn}
  \hspace{1.2em}
  \href{https://scholar.google.com/citations?user=clLJfm8AAAAJ}{\faGraduationCap~Google Scholar}
\end{center}

%----------- SUMMARY -----------
\section{Summary}
\vspace{-1pt}
{\small GPU-systems and ML-platform engineer focused on multi-node / multi-GPU training and inference: CUDA kernels and megakernels, distributed serving internals (vLLM), cross-rank GPU communication, and Kubernetes-managed B200 GPU clusters.}
\vspace{4pt}

%----------- EXPERIENCE -----------
\vspace{-2pt}
\section{Experience}
\resumeHeadingListStart
  \resumeSubheading{\href{https://www.microsoft.com/en-us/research/lab/microsoft-research-india/}{Microsoft Research}}{07/2024 -- Present}{Research Fellow}{}
    \resumeItemListStart
      \resumeItemPlain{Managed multi-node distributed training workloads on Kubernetes across Microsoft's NVIDIA B200 GPU clusters: pod scheduling, GPU resource allocation, and distributed-throughput optimization for large-scale code-model training.}
      \resumeItemPlain{Engineered a fused multi-GPU inference megakernel for DeepSeek-V2 (236B) on NVIDIA B200, merging all compute with tensor/expert-parallel communication into one kernel; handled cross-rank exchange via multimem PTX multicast (direct peer-to-peer NVLink writes), with block-scaled FP8, warp specialization, and mbarrier-based synchronization.}
      \resumeItemPlain{Trained an LLM on par with Gemini and GPT-4o for code generation, instruction-following, and reasoning via post-pretraining and online/offline RL.}
      \resumeItemPlain{Built the world's first model and pipeline for end-to-end autoformalization and verification of Rust programs; the GRPO-based RL recipe (with multiplicative gating rewards) lifted the verified-program rate from 1\% to 47\% over the SFT baseline; our 4B Verus-LM outperforms Kimi-K2.5 (1T) on competitive benchmarks.}
      \resumeItemPlain{Developing foundational algorithms that exploit reinforcement-learning gradient-optimization dynamics to improve representation efficiency, stability, and generalization.}
    \resumeItemListEnd

  \resumeSubheading{\href{https://numpy.org/}{NumPy (numpy-quaddtype)}}{12/2024 -- Present}{Maintainer}{}
    \resumeItemListStart
      \resumeItemPlain{Authored and maintain a cross-platform, true IEEE-compliant 128-bit quadruple-precision floating-point data type for NumPy, along with corresponding efficient and performant BLAS routines.}
      \resumeItemPlain{Review and merge core PRs affecting dtype architecture, numerical semantics, and cross-platform behavior; focus on dtype infrastructure, numerical correctness, and long-term ABI stability.}
      \resumeItemPlain{Work with contributors across architectures on portability, performance, and ABI concerns.}
    \resumeItemListEnd

  \resumeSubheading{\href{https://quansight.com/}{Quansight}}{07/2024 -- 09/2024}{Open Source Intern -- Numerical Types \& Precision (NumPy OSS)}{}
    \resumeItemListStart
      \resumeItemPlain{Built multi-platform distribution and testing infrastructure handling compiler toolchains, SIMD availability, FMA support, and architecture-specific behavior.}
      \resumeItemPlain{Led development of NumPy-QuadDType (cross-platform 128-bit float dtype); implemented casting, ufunc dispatch, scalar types, memory model, and serialization.}
    \resumeItemListEnd

  \resumeSubheading{dataX.ai}{05/2022 -- 10/2022}{Machine Learning Engineer Intern}{}
    \resumeItemListStart
      \resumeItemPlain{Deployed vision/language models via NVIDIA Triton Inference Server (ONNX conversion pipeline), cutting VM load by 12\%.}
      \resumeItemPlain{Wrote custom CUDA kernels for 3D medical-scan processing, achieving a 2x speedup in segmentation over existing GPU solutions.}
    \resumeItemListEnd
\resumeHeadingListEnd

%----------- PROJECTS -----------
\vspace{-2pt}
\section{Projects}
\resumeItemListStart
  \resumeItem{Free-Threaded vLLM Serving (Python no-GIL, CUDA)}{Re-architected vLLM's multi-process inference server into a single free-threaded (no-GIL) process, replacing ZMQ/msgpack and SHM/pickle IPC between the API layer, scheduler, and per-GPU workers with in-process reference passing, and removing the single-uvloop API-server bottleneck via per-frontend threads. Cut inter-engine coordination overhead $\sim$13.3x (94.7 ms to 7.1 ms wall-clock union) and P99 inter-token latency from $\sim$540 ms to $\sim$112 ms; best configuration reached 61.96 req/s vs 59.33 stock. Made CUDA graphs compatible in the decode path and added a thread-local tokenizer.}
  \resumeItem{Bare-Bones Inference (CUDA, C++)}{A minimal inference stack with engine-grade features such as batching and scheduling, but serving a single fused megakernel instead of a conventional multi-kernel pipeline. Built DeepSeek-V2 (236B)'s block-scaled FP8-quantized, multi-GPU megakernel for NVIDIA B200 (sm\_100) GPUs that fuses all computation and TP/EP communication into one kernel, handling cross-rank exchange via \texttt{multimem} PTX multicast writes (directly writing into peer ranks' memory at once). Achieves 13.11 ms/token decode latency, beating vLLM's optimized serving latency at $B{=}1$ by 2x and eager-mode by ${>}10$x.}
  \resumeItem{Blackwell Flash-Attention Kernels (CUDA, PTX)}{Prototyped flash-attention kernels for NVIDIA Blackwell (B200, sm\_100) using tcgen05 2SM MMA and warp specialization, with mbarrier-based producer/consumer synchronization and a persistent-CTA design; profiled against roofline limits using CuTe layouts.}
  \resumeItem{NVFP4 MoE Training System (CUDA, C++)}{Built a Mixture-of-Experts training system from scratch targeting NVIDIA B200 with NVFP4 (4-bit) numerics, focused on low-precision training stability and multi-GPU throughput.}
  \resumeItem{\href{https://github.com/numpy/numpy-quaddtype}{Numpy-QuadDType} (C, C++, Python)}{A cross-platform 128-bit (quadruple-precision) floating-point dtype for NumPy with 100k+ downloads. A portable alternative to long double with full casting, ufunc dispatch, scalar types, and serialization that behaves consistently across compilers and architectures.}
  \resumeItem{\href{https://github.com/SwayamInSync/QBLAS}{QBLAS} (C, C++)}{A high-performance BLAS library for IEEE-754 binary128 (quadruple) precision, implementing optimized linear-algebra kernels that bring quad-precision numerics to workloads unsupported by standard double-precision libraries.}
  \resumeItem{\href{https://github.com/SwayamInSync/cpp-verify}{cpp-verify} (C++, LLVM, SMT)}{Extends C++ with first-class formal-verification constructs (pre/post-conditions and invariants), lowering specifications through an LLVM-based pipeline and discharging the resulting proof obligations to SMT solvers.}
  \resumeItem{\href{https://github.com/swayamInSync/mira}{MIRA -- Multimodal Image Reconstruction with Attention} (PyTorch)}{A transformer-based architecture for single-view text/image-to-3D reconstruction using ViT encoders and triplane decoders with cross-attention. Generates 3D mesh and video in under 10s on an A100, with SDXL-integrated diffusion for end-to-end scene synthesis.}
  \resumeItem{\href{https://github.com/SwayamInSync/clothes-virtual-try-on}{Virtual Clothing Assistant} (PyTorch)}{End-to-end virtual try-on system: ResNet101 and UNet for garment/body segmentation, OpenPose for pose estimation, and a PyTorch VITON pipeline for realistic garment warping. Reaches SSIM 0.895 with 500+ GitHub stars.}
\resumeItemListEnd

%----------- PUBLICATIONS -----------
\vspace{-2pt}
\section{Publications}
\resumeItemListStart
  \resumeItemPlain{\textbf{\href{https://www.microsoft.com/en-us/research/publication/nextcoder-robust-adaptation-of-code-lms-to-diverse-code-edits/}{NextCoder: Robust Adaptation of Code LMs to Diverse Code Edits}} (ICML 2025; ICLR 2025 DL4Code). \textbf{Swayam Singh}, Tushar Aggarwal, et al. Synthetic-data pipeline + SeleKT adaptation; 32B matches Gemini and beats GPT-4o on harder edit tasks, purely via SFT.}
  \resumeItemPlain{\textbf{\href{https://arxiv.org/abs/2407.03941}{Narrow Transformer: StarCoder-Based Java-LM For Desktop}} (arXiv:2407.03941, 2024). Kamalkumar Rathinasamy, \ldots, \textbf{Swayam Singh}, et al. A compact, Java-specialized code LM for desktop hardware.}
  \resumeItemPlain{\textbf{\href{https://arxiv.org/abs/2308.07124}{OctoPack: Instruction Tuning Code Large Language Models}} (ICLR 2024 Spotlight; NeurIPS 2023 Instruction Workshop). Niklas Muennighoff, \ldots, \textbf{Swayam Singh}, et al. Instruction tuning of code models via natural-language Git commits (CommitPack).}
  \resumeItemPlain{\textbf{\href{https://arxiv.org/abs/2305.06161}{StarCoder: May the Source Be With You!}} (TMLR 2023). Raymond Li, \ldots, \textbf{Swayam Singh}, et al.\ (BigCode). A 15.5B open code LLM trained on 1T tokens of permissively licensed code.}
\resumeItemListEnd

%----------- SKILLS -----------
\vspace{-2pt}
\section{Skills Summary}
\resumeItemListStart
  \resumeItem{Systems, GPU \& Serving}{CUDA kernels and megakernels, multi-GPU / multi-node deployment, tensor/expert parallelism, cross-rank collectives (multimem / NVLink), CUDA graphs, FP8 quantization, inference serving (vLLM internals), NVIDIA Triton Inference Server, SIMD/FMA, performance optimization}
  \resumeItem{Infra \& Distributed}{Kubernetes (multi-node GPU workloads, scheduling, resource allocation), distributed training, cross-platform C-extension packaging \& CI}
  \resumeItem{Languages}{C, C++, CUDA, CuTeDSL, Python (incl.\ free-threaded / no-GIL)}
  \resumeItem{ML \& Training}{LLM post-training, reinforcement learning (online/offline RL), supervised fine-tuning, code-generation models}
  \resumeItem{Tools \& Frameworks}{PyTorch, CuPy, NumPy (C-API internals), LLVM, SMT solvers}
  \resumeItem{Open Source}{Core NumPy maintenance, cross-platform C-extension packaging, code review}
\resumeItemListEnd

%----------- HONORS & AWARDS -----------
\vspace{-2pt}
\section{Honors and Awards}
\resumeItemListStart
  \resumeItemPlain{\textbf{2024} -- \textbf{Kaggle Competition Expert}: Bronze medal (top 7\%) in UBC-OCEAN; top 3\% in the \textit{30 Days of ML} challenge.}
  \resumeItemPlain{\textbf{2024} -- Invited to \textbf{Google Research Week} (keynote by Jeff Dean).}
  \resumeItemPlain{\textbf{2024} -- \textbf{OctoPack} accepted as a \textbf{Spotlight (top 5\%)} at ICLR 2024.}
  \resumeItemPlain{\textbf{2023} -- Selected for the \textbf{Amazon ML Summer School 2023}.}
  \resumeItemPlain{\textbf{2023} -- \textbf{Clothes Virtual Try-On} crossed 500+ GitHub stars.}
\resumeItemListEnd

%----------- INVITED TALKS -----------
\vspace{-2pt}
\section{Invited Talks}
\resumeItemListStart
  \resumeItemPlain{\textbf{MAMBA: Zero to Hero}, invited talk on State Space Models at Cohere for AI.}
  \resumeItemPlain{\textbf{Provably-Correct Code and Efficient Sparse Training of LLMs}, internal research talk at Microsoft Research.}
  \resumeItemPlain{\textbf{Foundations of Machine Learning}, a GDG On-Campus session on the ML landscape, tooling ecosystem, and emerging directions.}
  \resumeItemPlain{\textbf{From Deep Learning to Large Language Models}, a talk on the foundations of modern AI: deep learning, generative models, and LLMs.}
\resumeItemListEnd

\end{document}
